\section{Local divisors and Schubert charts}\label{sec:algebra}

All schemes are Noetherian, and varieties are over $\C$ unless a field extension is specified. On a normal integral scheme, Weil divisor classes correspond to rank-one reflexive sheaves. A class vanishes in a local class group precisely when the corresponding module is free of rank one. We use the divisor--sheaf correspondence, the height-one criterion for reflexivity, and Serre's normality criterion \cite[Chapter~II, Section~6]{Har77}\cite[Tags 0EBJ, 0EBM, 0AVB, 031S]{The26}.

\begin{lemma}[Localization of divisor classes]\label{alg:localization}
Let $R$ be a normal Noetherian domain, let $0\ne t\in R$, and let $E_1,\ldots,E_s$ be the height-one primes containing $t$. Then
\[
 \bigoplus_{i=1}^s\Z[E_i]\longrightarrow\Cl(R)
 \longrightarrow\Cl(R[t^{-1}])\longrightarrow0
\]
is exact. Polynomial and Laurent polynomial extensions have the same class group as $R$.
\end{lemma}
\begin{proof}
Every prime divisor of $D(t)$ has a unique prime-divisor closure in $\Spec R$. If a divisor restricts to a principal divisor, subtract the divisor of the same rational function on $\Spec R$; the difference is supported on $V(t)$. This proves exactness, including the required kernel generation. For polynomial invariance, divisors dominating $\Spec R$ become principal over $\Frac(R)[T]$, while vertical valuations are extensions of the height-one valuations of $R$; Gauss's lemma identifies the relations. Localizing in $T$ gives Laurent invariance \cite[Chapter~II, Propositions~6.5--6.6]{Har77}\cite{Fos73}.
\end{proof}

The classes $[E_i]$ need not be linearly independent in $\Cl(R)$. The same closure argument applies to any nonempty open subset $U$ of a normal integral variety $X$: $\Cl(U)$ is the quotient of $\Cl(X)$ by the subgroup generated by prime divisors in $X\setminus U$. In particular it applies to the affine big cell of a Schubert variety even though its complement need not be the divisor of a single global function.

\begin{lemma}[Flat base change and reflection of invertibility]\label{alg:flat}
Let $A\to B$ be a flat homomorphism of Noetherian rings and let $M$ be a finite $A$-module. There are natural isomorphisms
\[
 \begin{aligned}
 \Hom_A(M,A)\otimes_A B&\simeq\Hom_B(M\otimes_A B,B),\\
 M^{**}\otimes_A B&\simeq(M\otimes_A B)^{**}.
 \end{aligned}
\]
For normal domains, flat base change therefore preserves finite rank-one reflexive modules and their reflexive tensor powers. Invertibility descends under faithful flatness; over a local source, freeness of rank one descends.
\end{lemma}
\begin{proof}
A finite presentation expresses the dual as a kernel between finite free modules. Flat tensor product preserves this kernel \cite[Tag 0C6I]{The26}. The dual is again finite over a Noetherian ring, so applying the argument twice proves the double-dual formula, and applying it to tensor powers proves the corresponding assertion. A rank-one module remains rank one after the indicated flat extension of domains. The last assertion is faithfully flat descent for finite locally free modules \cite[Tag 05B2]{The26}.
\end{proof}

\begin{lemma}[Generic smooth factors]\label{alg:generic}
Let $(R,\mathfrak m)$ be a normal Noetherian local domain and set $R(T)=R[T]_{\mathfrak mR[T]}$. A finite rank-one reflexive $R$-module $M$ is free if and only if $M\otimes_RR(T)$ is free. Consequently passage to the generic point of an affine cell factor reflects the Cartier condition. The analogous multivariable localization is denoted by $R(T_1,\ldots,T_a)$, and the same assertion holds for it.
\end{lemma}
\begin{proof}
This is a faithfully flat local extension with residue field extension $\kappa(\mathfrak m)\subset\kappa(\mathfrak m)(T)$. The minimal number of generators is unchanged. If the extended module is free of rank one, $M$ is cyclic, and a cyclic torsion-free module of rank one over a domain is free. The forward implication is immediate.
\end{proof}

\subsection{Charts and conventions}\label{sec:charts}
Write $W$ for the Weyl group, $C_u=BuB/B$ for an ordinary cell, and $C^u=B^-uB/B$ for an opposite cell. For $u\le w$, the Kazhdan--Lusztig slice $N_{u,w}=X_w\cap C^u$ is affine, has distinguished point $u$, and has dimension $\ell(w)-\ell(u)$. A standard open neighbourhood $\Omega_u$ gives an actual scheme isomorphism
\[
 X_w\cap\Omega_u\simeq C_u\times N_{u,w},
\]
compatible with Schubert strata \cite[Lemma~3.2]{WY08}\cite[Lemma~2.1]{KWY13}. Normality and Cohen--Macaulayness of the slices follow from this product and the corresponding results for Schubert varieties \cite[Theorems~2.1.1, 2.2.3 and Proposition~2.2.5]{Bri05}. The opposite cell supplies a contracting positive grading with degree-zero part $\C$.

When identifying prime ideals, we use scheme-theoretic reducedness of intersections of Schubert varieties and their opposite analogues \cite{Ram87}\cite[Section~2]{Kum17}. A calculation of root tangent directions by itself would not establish the required ideal equality.

For a global Weil divisor $D$, its non-Cartier locus $\NC(D)$ is the closed nonfree locus of the coherent reflexive sheaf $\OO_X(D)$. If $D$ is $B$-stable, so is this locus; its irreducible components are then Schubert subvarieties, because there are finitely many $B$-orbits. Any local Weil class extends to a global Weil class by closing its prime divisors. It can also be represented by a $B$-stable boundary divisor: the big cell is affine space, its class group is zero, and Lemma~\ref{alg:localization} supplies the representation.
